\section{The calculus \texorpdfstring{$\CICm$}{CIC-}}
\label{sec:calculus}

We work with a calculus $\CICm$, a predicative fragment of the
Calculus of Inductive Constructions with an impredicative universe of
propositions, primitive logical connectives, and parametric inductive
data types.  The design follows the intermediate representation
$\mathrm{CIC}_0$ of~\cite{CzajkaKaliszyk2018} in having an explicit
syntactic distinction between propositions and data, and primitive
connectives; it deviates from full CIC in ways that we record
explicitly in \cref{rem:fragment} and discuss throughout.  The two
restrictions that carry mathematical weight are the absence of
proof-dependent data (\cref{lem:proof-confinement}) and the
restriction of inductive types to parametric data types.

\subsection{Syntax}

\begin{definition}[Sorts and terms]
\label{def:terms}
The set of \emph{sorts} is $\Sorts = \{\SProp, \STypeI, \STypeII\}$.
The set of \emph{terms} is generated by the grammar
\begin{align*}
t,u,A,B,\varphi,\psi \Coloneqq{}
  & x \mid c \mid s \mid \Pi x{:}A.\,B \mid \lambda x{:}A.\,t \mid t\,u
    \mid \case{I,\tup{\tau},Q}{t}{\tup{u}} \\
  & \mid \top \mid \bot \mid \varphi \wedge \psi \mid \varphi \vee \psi
    \mid \varphi \leftrightarrow \psi \mid t \ceq_A u \mid \cexists{x}{A}\varphi
\end{align*}
where $x$ ranges over variables, $c$ over constants and $s$ over
sorts.  We write $A \to B$ for $\Pi x{:}A.\,B$ when $x \notin \FV(B)$,
and $\neg\varphi$ for $\varphi \to \bot$.  Universal quantification is
not a separate former: $\forall x{:}A.\,\varphi$ \emph{is} the
dependent product $\Pi x{:}A.\,\varphi$.  Terms are identified up to
$\alpha$-conversion; $\FV(t)$ denotes the set of free variables;
$t[u/x]$ denotes capture-avoiding substitution, and $t[\tup{u}/\tup{x}]$
its simultaneous variant.
\end{definition}

\begin{definition}[Environments]
\label{def:env}
A \emph{global environment} $\Env$ is a finite sequence of
declarations of the following forms:
\begin{enumerate}[label=(\roman*)]
\item a \emph{definition} $c \coloneqq t : T$;
\item a \emph{parameter declaration} $c : T$ (a constant without a body;
  premises available to the hammer are of this form once their proof
  terms are discarded);
\item an \emph{inductive declaration}
  \[
    \Ind\bigl(I;\; A_1{:}\STypeI, \ldots, A_p{:}\STypeI;\;
        c_1 : \Pi\tup{A}{:}\STypeI.\,\sigma_{1,1} \to \cdots \to \sigma_{1,n_1} \to I\,\tup{A},\;
        \ldots,\;
        c_k : \cdots\bigr)
  \]
  introducing a \emph{data inductive type} $I$ with $p \geq 0$
  parameters and $k \geq 0$ constructors $c_1, \ldots, c_k$.  Each
  constructor argument type $\sigma_{j,l}$ is required to be either
  (a) a term with $A_1{:}\STypeI, \ldots, A_p{:}\STypeI \vdash
  \sigma_{j,l} : \STypeI$ in which $I$ does not occur, or (b) exactly
  the term $I\,\tup{A}$ (a \emph{recursive argument}).
\end{enumerate}
A \emph{context} $\Gamma$ is a finite sequence $x_1{:}A_1, \ldots,
x_n{:}A_n$ of variable declarations.  The environment is a global
parameter of all judgements below; we keep it implicit and write
$\Gamma \vdash t : T$ for the typing judgement.
\end{definition}

The inductive scheme of \cref{def:env}(iii) covers the parametric
data types relevant to monomorphisation --- $\mathsf{nat}$,
$\mathsf{bool}$, $\mathsf{unit}$, $\mathsf{list}\,A$,
$\mathsf{prod}\,A\,B$, $\mathsf{option}\,A$, binary trees, and so on
--- while excluding indexed families and nested or higher-order
recursive occurrences.  Constructor arguments do not depend on one
another.  Inductively defined \emph{propositions}
($\mathsf{Forall}$, $\mathsf{le}$, \ldots) are not part of the
calculus: an inductive predicate enters a problem as a parameter
constant $c : \tup{\sigma} \to \SProp$ together with premises
axiomatising it --- for instance its introduction rules and its
inversion principle, all of which are provable propositions of the
ambient proof assistant.  This is not a loss of generality for our
purposes: the soundness of any axiom that is the translation of a
\emph{provable} proposition follows from the premise-truth part of
\cref{thm:base-soundness} below, so inductive predicates require no
dedicated treatment.  Only \emph{data} inductive types need built-in
support, because the translation emits genuinely new axioms for them
(injectivity, discrimination, inversion), and precisely those axioms
are the subject of the guard-erasure analysis of
\cref{sec:erasure}.

\subsection{Reduction and conversion}

\begin{definition}[Reduction]
\label{def:reduction}
The relation $\tored$ is the contextual closure of:
\begin{itemize}
\item ($\beta$)\; $(\lambda x{:}A.\,t)\,u \tored t[u/x]$;
\item ($\iota$)\; $\case{I,\tup{\tau},Q}{c_j\,\tup{\tau}\,\tup{v}}{\tup{u}}
      \tored u_j\,\tup{v}$, where $c_j$ is the $j$-th constructor of $I$;
\item ($\delta$)\; $c \tored t$, for every definition
      $c \coloneqq t : T$ in the environment.
\end{itemize}
$\toredstar$ is its reflexive--transitive closure and $\conv$
(\emph{conversion}) the generated equivalence.  We write $t \Downarrow
u$ if $t \toredstar u$ and $u$ is in normal form.
\end{definition}

\subsection{Typing}

\begin{definition}[Typing]
\label{def:typing}
The judgement $\Gamma \vdash t : T$ is generated by the following
rules.  The PTS core uses the axioms $\mathcal{A} = \{\SProp :
\STypeI,\; \STypeI : \STypeII\}$ and the rule set
\[
  \mathcal{R} \;=\;
  \setof{(s, \SProp, \SProp)}{s \in \Sorts}
  \;\cup\;
  \setof{(\mathsf{Type}_i, \mathsf{Type}_j, \mathsf{Type}_{\max(i,j)})}{i,j \in \{1,2\}} .
\]
\begin{gather*}
\rulname{ax}\rul{(s_1 : s_2) \in \mathcal{A}}{\Gamma \vdash s_1 : s_2}
\qquad
\rulname{var}\rul{(x{:}A) \in \Gamma \quad \Gamma \vdash A : s}{\Gamma \vdash x : A}
\qquad
\rulname{const}\rul{(c : T) \text{ or } (c \coloneqq t : T) \text{ in } \Env}{\Gamma \vdash c : T}
\\[1ex]
\rulname{prod}\rul{\Gamma \vdash A : s_1 \quad \Gamma, x{:}A \vdash B : s_2 \quad (s_1,s_2,s_3) \in \mathcal{R}}{\Gamma \vdash \Pi x{:}A.\,B : s_3}
\\[1ex]
\rulname{lam}\rul{\Gamma, x{:}A \vdash t : B \quad \Gamma \vdash \Pi x{:}A.\,B : s}{\Gamma \vdash \lambda x{:}A.\,t : \Pi x{:}A.\,B}
\\[1ex]
\rulname{app}\rul{\Gamma \vdash t : \Pi x{:}A.\,B \quad \Gamma \vdash u : A}{\Gamma \vdash t\,u : B[u/x]}
\qquad
\rulname{conv}\rul{\Gamma \vdash t : A \quad \Gamma \vdash B : s \quad A \conv B}{\Gamma \vdash t : B}
\end{gather*}
For an inductive declaration as in \cref{def:env}(iii), with
$\tup{\tau}$ a vector of $p$ terms:
\begin{gather*}
\rulname{ind}\rul{}{\Gamma \vdash I : \underbrace{\STypeI \to \cdots \to \STypeI}_{p} \to \STypeI}
\\[1ex]
\rulname{constr}\rul{}{\Gamma \vdash c_j : \Pi \tup{A}{:}\STypeI.\,\sigma_{j,1} \to \cdots \to \sigma_{j,n_j} \to I\,\tup{A}}
\\[1ex]
\rulname{case}\rul{\begin{array}{c}
  \Gamma \vdash t : I\,\tup{\tau} \qquad
  \Gamma \vdash Q : I\,\tup{\tau} \to s \qquad s \in \{\SProp, \STypeI\} \\
  \Gamma \vdash u_j : \sigma_{j,1}[\tup{\tau}/\tup{A}] \to \cdots \to
     \sigma_{j,n_j}[\tup{\tau}/\tup{A}] \to Q\,(c_j\,\tup{\tau}\,\tup{y})
     \;\text{ for } j = 1, \ldots, k
  \end{array}}
  {\Gamma \vdash \case{I,\tup{\tau},Q}{t}{\tup{u}} : Q\,t}
\end{gather*}
(the branch typing is understood as: $u_j$ has the displayed closed
product type, in which $\tup{y}$ are the bound argument variables).
Finally, the logical formers:
\begin{gather*}
\rul{}{\Gamma \vdash \top : \SProp} \qquad
\rul{}{\Gamma \vdash \bot : \SProp} \qquad
\rul{\Gamma \vdash \varphi : \SProp \quad \Gamma \vdash \psi : \SProp \quad \circ \in \{\wedge, \vee, \leftrightarrow\}}{\Gamma \vdash \varphi \circ \psi : \SProp}
\\[1ex]
\rul{\Gamma \vdash A : \STypeI \quad \Gamma \vdash t : A \quad \Gamma \vdash u : A}{\Gamma \vdash t \ceq_A u : \SProp}
\qquad
\rul{\Gamma \vdash A : \STypeI \quad \Gamma, x{:}A \vdash \varphi : \SProp}{\Gamma \vdash \cexists{x}{A}\varphi : \SProp}
\end{gather*}
Well-formedness of contexts and environments is defined as usual: each
declared type must be typable by a sort (for definitions, the body
must have the declared type), and the constructor argument condition
of \cref{def:env}(iii) must hold.  We consider only well-formed
environments and contexts from now on.
\end{definition}

Note which products exist.  Since $\mathcal{R}$ contains no rule of
the form $(\SProp, \mathsf{Type}_i, -)$, there are \emph{no functions
from proofs to data or to types}: a product $\Pi x{:}\varphi.\,B$ with
$\varphi : \SProp$ is well-formed only when $B : \SProp$.
Quantification over propositions ($\Pi A{:}\SProp.\,\varphi$, sort of
the domain $\STypeI$) and over data types ($\Pi
A{:}\STypeI.\,\varphi$, sort of the domain $\STypeII$) into $\SProp$
are both available, giving the impredicativity of $\SProp$ familiar
from CIC.  Polymorphic data such as $\Pi A{:}\STypeI.\,A \to
\mathsf{list}\,A$ lives in $\STypeII$; type constructors such as
$\mathsf{list} : \STypeI \to \STypeI$ have kinds in $\STypeII$.  There
is no cumulativity: $\SProp$, $\STypeI$ and $\STypeII$ are unrelated
by subtyping.

\begin{definition}[Classification]
\label{def:classification}
Let $\Gamma \vdash t : T$.  We call $t$:
a \emph{proposition} if $T \conv \SProp$\,%
\footnote{Here and below, statements such as ``$T \conv \SProp$'' are
justified by uniqueness of types (\cref{prop:metatheory}): the
alternative types of $t$ agree up to conversion, and distinct sorts
are distinct normal forms, so the classification does not depend on
the choice of $T$.};
a \emph{proof} if $\Gamma \vdash T : \SProp$;
a \emph{data type} if $T \conv \STypeI$;
a \emph{data object} if $\Gamma \vdash T : \STypeI$;
and \emph{kind-level} otherwise (this covers $\STypeI$ itself, kinds
such as $\STypeI \to \STypeI$, type constructors, polymorphic
functions, and $\STypeII$-sorted products generally).
A variable $x$ with $(x{:}A) \in \Gamma$ is a \emph{proof variable}
if $A$ is a proposition, a \emph{data variable} if $A$ is a data
type, and a \emph{type variable} if $A = \STypeI$.
\end{definition}

\subsection{Metatheory}

The following standard properties are inherited from the general
theory of Pure Type Systems and of CIC; we state them for reference
and use them freely.  $\CICm$ is a functional (singly sorted) PTS
extended with parametric inductive types satisfying the usual
strict-positivity condition, so the results below are instances of
the literature
\cite{Barendregt1992,Geuvers1993,Werner1994,PaulinMohring1993}.

\begin{proposition}[Standard metatheory]
\label{prop:metatheory}
\leavevmode
\begin{enumerate}[label=(\arabic*)]
\item \emph{(Confluence)} $\tored$ is confluent on well-typed terms.
\item \emph{(Substitution)} If $\Gamma, x{:}A, \Delta \vdash t : T$
  and $\Gamma \vdash u : A$, then
  $\Gamma, \Delta[u/x] \vdash t[u/x] : T[u/x]$.
\item \emph{(Subject reduction)} If $\Gamma \vdash t : T$ and
  $t \tored t'$ then $\Gamma \vdash t' : T$.
\item \emph{(Uniqueness of types)} If $\Gamma \vdash t : T_1$ and
  $\Gamma \vdash t : T_2$ then $T_1 \conv T_2$.
\item \emph{(Strong normalisation)} Every well-typed term is strongly
  $\beta\iota\delta$-normalising.
\item \emph{(Sort classification)} If $\Gamma \vdash t : T$, then
  either $T \conv \STypeII$, or there is a unique sort
  $s \in \Sorts$ with $\Gamma \vdash T : s$.  Consequently the five
  classes of \cref{def:classification} are mutually exclusive and
  exhaustive on well-typed terms.
\end{enumerate}
\end{proposition}

The next two lemmas are specific to $\CICm$ and carry real weight in
this paper, so we prove them in full.  The first says that proofs are
hermetically confined: a term that is not itself a proof contains no
proofs anywhere inside it.  This is what makes a \emph{total},
proof-free translation possible on the fragment
(\cref{sec:translation}), and it plays the role that proof
irrelevance plays in the piPTS soundness proof
of~\cite{Czajka2016}: the known counterexample to soundness of
shallow embeddings without proof irrelevance
(\cite[Example~52]{Czajka2016}: a predicate $q : p \to \SProp$
applied to two different proof variables) is not expressible in
$\CICm$, because the type $p \to \SProp$ with $p : \SProp$ requires
the excluded rule $(\SProp, \STypeII, \cdot)$.

\begin{lemma}[Proof confinement]
\label{lem:proof-confinement}
Let $\Gamma \vdash t : T$ and suppose $t$ is not a proof.  Then no
subterm occurrence of $t$ (including type annotations of binders) is
a proof relative to its local context, and in particular no free or
bound variable of $t$ is a proof variable.
\end{lemma}

\begin{proof}
By induction on the structure of $t$, using the generation
(inversion) properties of the typing rules.  We check that every
immediate subterm of $t$, in its local context, is again a non-proof;
the statement then follows by the induction hypothesis applied to the
immediate subterms.

\emph{Case $t$ a variable, constant or sort:} immediate (no proper
subterms).

\emph{Case $t = u\,v$:} by generation, $\Gamma \vdash u : \Pi
x{:}A.\,B$ with $\Gamma \vdash v : A$ and $T \conv B[v/x]$.  First,
$u$ is not a proof: the type $\Pi x{:}A.\,B$ of $u$ has some sort
$s_3$ with $(s_1, s_2, s_3) \in \mathcal{R}$, and if $s_3 = \SProp$
then, by the shape of $\mathcal{R}$, $s_2 = \SProp$; then $B :
\SProp$, hence $B[v/x] : \SProp$ by
\cref{prop:metatheory}(2), so $t$ would be a proof, contradiction.
So $s_3 \neq \SProp$ and $u$ is not a proof.  Second, $v$ is not a
proof: if $A$ were a proposition, i.e.\ $A : \SProp$, then the
product rule used for $\Pi x{:}A.\,B$ has $s_1 = \SProp$, and the
only rules in $\mathcal{R}$ with first component $\SProp$ are
$(\SProp, \SProp, \SProp)$; thus $B : \SProp$ and again $t$ is a
proof, contradiction.

\emph{Case $t = \lambda x{:}A.\,u$:} by generation $T \conv \Pi
x{:}A.\,B$ with $\Gamma, x{:}A \vdash u : B$, and $\Pi x{:}A.\,B : s_3$
with $s_3 \neq \SProp$ since $t$ is not a proof.  As in the previous
case, $s_3 \neq \SProp$ forces $s_1 \neq \SProp$ (all rules with
$s_1 = \SProp$ have $s_3 = \SProp$), so $A$ is not a proposition and
$x$ is not a proof variable; moreover $B$'s sort $s_2$ satisfies
$(s_1, s_2, s_3) \in \mathcal{R}$ with $s_3 \neq \SProp$, hence $s_2
\neq \SProp$, so $u$ is not a proof in $\Gamma, x{:}A$.  The
annotation $A$ has type $s_1 \neq \SProp$, so $A$ is not a proof
either.

\emph{Case $t = \Pi x{:}A.\,B$:} $t$ has a sort as its type, so $t$
is never a proof, and its immediate subterms $A$ and $B$ have sorts
as types, so neither is a proof.

\emph{Case $t = \case{I,\tup{\tau},Q}{u}{\tup{u}}$:} the scrutinee
satisfies $\Gamma \vdash u : I\,\tup{\tau} : \STypeI$, so $u$ is a
data object, not a proof.  The parameters $\tup{\tau}$ are data types
($\STypeI$-typed), not proofs.  The motive $Q$ has type
$I\,\tup{\tau} \to s$, a kind or a $\STypeII$-sorted product, so $Q$
is not a proof.  The branches $u_j$ have $\STypeII$- or
$\STypeI$-sorted product types when $s = \STypeI$; when $s = \SProp$
they have types of sort $\SProp$ --- but then $Q\,u : \SProp$ and $t$
itself is a proof, contradicting the hypothesis.  So when $t$ is not
a proof, $s = \STypeI$ and no immediate subterm is a proof.

\emph{Case $t$ a connective, equation or existential:} $t : \SProp$,
so $t$ is a proposition, not a proof.  Immediate subterms:
subformulas are propositions; the sides of $t_1 \ceq_A t_2$ are data
objects and $A$ a data type; the domain of $\cexists{x}{A}\varphi$ is
a data type and the body a proposition.  None are proofs.

For the ``in particular'' part: a free proof variable would itself be
a proof subterm occurrence; a bound proof variable requires a binder
$\lambda x{:}\varphi$ or $\Pi x{:}\varphi$ or $\cexists{x}{\varphi}$
with $\varphi : \SProp$ inside $t$.  The $\lambda$-case was excluded
above whenever the $\lambda$ is not a proof; a subterm $\lambda$ that
\emph{is} a proof cannot occur, since we showed no subterm of $t$ is
a proof.  A product $\Pi x{:}\varphi.\,\psi$ with a proposition
domain forces $\psi : \SProp$; such a subterm is an implication
$\varphi \to \psi$: its bound variable cannot occur in $\psi$,
because $x \in \FV(\psi)$ with $x$ a proof variable would make some
subterm of $\psi$ contain the free proof variable $x$, and $\psi$ is
a proposition, hence not a proof, so by the induction hypothesis
$\psi$ has no free proof variables.  The existential former binds
only data variables by its typing rule.
\end{proof}

\begin{corollary}
\label{cor:impl-nondep}
If $\Gamma \vdash \Pi x{:}\varphi.\,\psi : \SProp$ with $\Gamma
\vdash \varphi : \SProp$, then $x \notin \FV(\psi)$.  Products over
propositions are plain implications.
\end{corollary}

The second lemma states that instantiating a type variable by a
closed data type leaves the classification of every subterm
unchanged.  It is the reason why the type-directed translation of
\cref{sec:translation} interacts predictably with the
monomorphisation of \cref{sec:mono}.

\begin{lemma}[Sort invariance under type instantiation]
\label{lem:sort-invariance}
Let $\Gamma, A{:}\STypeI, \Delta \vdash t : T$ and let $\tau$ be a
term with $\Gamma \vdash \tau : \STypeI$.  Then
$\Gamma, \Delta[\tau/A] \vdash t[\tau/A] : T[\tau/A]$, and $t[\tau/A]$
belongs to the same class of \cref{def:classification} (proposition,
proof, data type, data object, kind-level) as $t$.
\end{lemma}

\begin{proof}
The typing statement is \cref{prop:metatheory}(2).  For the
classification, first note that sorts are closed terms, so $s[\tau/A]
= s$ for every sort $s$.

If $t$ is kind-level with $T \conv \STypeII$, then $T[\tau/A] \conv
\STypeII[\tau/A] = \STypeII$ (conversion is closed under
substitution), so $t[\tau/A]$ is kind-level.  Otherwise, by
\cref{prop:metatheory}(6) there is a unique sort $s$ with $\Gamma,
A{:}\STypeI, \Delta \vdash T : s$.  By \cref{prop:metatheory}(2)
applied to this judgement, $\Gamma, \Delta[\tau/A] \vdash T[\tau/A] :
s[\tau/A] = s$.  By uniqueness of types
(\cref{prop:metatheory}(4),(6)), $s$ is \emph{the} sort of the type
of $t[\tau/A]$.  Hence: $t$ a proof (sort of $T$ is $\SProp$) iff
$t[\tau/A]$ a proof; $t$ a data object iff $t[\tau/A]$ a data
object; and similarly for kind-level terms whose type has sort
$\STypeII$.

For the remaining two classes, which are defined by the type itself
rather than by its sort: if $T \conv \SProp$ then $T[\tau/A] \conv
\SProp$, so propositions map to propositions; if $T \conv \STypeI$
then $T[\tau/A] \conv \STypeI$, so data types map to data types.
Conversely, no new memberships in these classes can appear: the
classes are mutually exclusive and exhaustive
(\cref{prop:metatheory}(6)), and we have just shown each class to be
preserved, so the induced map on classes is the identity.
\end{proof}

\begin{remark}[Relation to full CIC]
\label{rem:fragment}
Compared to the calculus underlying Coq/Rocq, $\CICm$ omits:
cumulativity and the full universe hierarchy (we keep two predicative
levels, which suffice to type polymorphic premises and their
instances; universe polymorphism and the constraints machinery are
orthogonal to our concerns --- note that the implemented translation
\emph{drops} universe constraints altogether
\cite[\S3]{CzajkaKaliszyk2018}, which is one of its known sources of
theoretical unsoundness, whereas $\CICm$ is honestly stratified);
proof-dependent data, i.e.\ the rules $(\SProp, \mathsf{Type}_i,
\mathsf{Type}_i)$, hence strong sums such as \texttt{sig} and
proof-carrying records; indexed inductive families and
nested/mutual/higher-order recursion; general \texttt{fix} (recursive
functions can be accommodated as parameter constants whose guarded
defining equations enter as premises --- the translation
of~\cite{CzajkaKaliszyk2018} in any case converts \texttt{fix} into
exactly such equations, and our soundness framework only requires the
equations to be \emph{valid}, cf.\ \cref{def:admissible}); and
$\eta$-conversion.  The implemented CoqHammer translation handles the
omitted features by opaque naming (``lifting''); nothing in the
monomorphisation and guard-erasure analysis below depends on them.
\end{remark}
